\documentclass[11pt,a4paper]{article}
\usepackage[T1]{fontenc}
\usepackage{lmodern,amsmath,amssymb,amsthm,microtype,booktabs}
\usepackage[margin=25mm]{geometry}
\usepackage[colorlinks=true,allcolors=blue]{hyperref}
\newtheorem{theorem}{Theorem}
\newtheorem{lemma}{Lemma}
\newcommand{\R}{\mathbb R}
\title{Diagonal concavity does not imply concavity for sums of elementary symmetric polynomials}
\author{Dev D. Goyal}
\date{26 October 2026}
\usepackage[numbers]{natbib}
\usepackage{xurl}
\begin{document}
\maketitle
\begin{abstract}
An explicit family disproves Lachaume's diagonal concavity conjecture in four variables with strictly positive coefficients. For the family $f_M=1+M\sigma_1+\sigma_2+M\sigma_3+\sigma_4$, where $M>0$, the fourth root is concave on the positive diagonal exactly when $M\geq1$. It is concave on the full positive orthant exactly when $M=1$. Two polynomial identities prove the failure of concavity. Concavity at $M=1$ follows from the arithmetic--geometric mean inequality. The Lean formalization proves both classifications with Mathlib's concavity predicate and elementary symmetric functions.
\end{abstract}

\section{The diagonal concavity conjecture}
For $x=(x_1,x_2,x_3,x_4)$, define the elementary symmetric polynomials by
\[
 \sigma_j(x)=\sum_{1\leq i_1<\cdots<i_j\leq4}x_{i_1}\cdots x_{i_j},
 \qquad \sigma_0(x)=1.
\]
Write $\Gamma=(0,\infty)^4$. A function $H$ is concave on a convex set if
\[
 H(\theta x+(1-\theta)y)\geq\theta H(x)+(1-\theta)H(y)
 \quad(0\leq\theta\leq1)
\]
for all $x,y$ in that set.
Lachaume's Conjecture 1 \cite{lachaume} asserts an equivalence between concavity on the positive orthant and concavity on its positive diagonal for
\[
 \left(\sum_{j=0}^p a_j\sigma_j\right)^{1/p},\qquad a_j\geq0.
\]
Here the positive diagonal consists of the vectors $(t,\ldots,t)$ with $t>0$.

\begin{theorem}\label{thm:classification}
For $M>0$, put
\[
 f_M=1+M\sigma_1+\sigma_2+M\sigma_3+\sigma_4,
 \qquad H_M=f_M^{1/4}.
\]
Then
\begin{align*}
 t\longmapsto H_M(t,t,t,t)\text{ is concave on }(0,\infty)
 &\quad\Longleftrightarrow\quad M\geq1,\\
 H_M\text{ is concave on }\Gamma
 &\quad\Longleftrightarrow\quad M=1.
\end{align*}
Thus each $M>1$ disproves Conjecture 1. In particular, $(a_0,a_1,a_2,a_3,a_4)=(1,2,1,2,1)$ is a counterexample with all coefficients positive.
\end{theorem}
Each term in $f_M$ is nonnegative on $\Gamma$, and the constant term is one. Thus the fourth root is defined and smooth throughout the domain. The polynomial has degree four, every coefficient is positive, and the counterexample lies in the interior of the domain.

\section{Two curvature identities}
For a positive twice differentiable function $q$,
\begin{equation}\label{eq:root}
 (q^{1/4})''=\frac{4qq''-3(q')^2}{16q^{7/4}}.
\end{equation}
A twice differentiable function on an open interval is concave exactly when its second derivative is nonpositive. This follows from the mean value theorem and the monotonicity of its first derivative.

The diagonal polynomial is
\[
 P_M(t)=f_M(t,t,t,t)=1+4Mt+6t^2+4Mt^3+t^4.
\]
Expansion gives
\begin{equation}\label{eq:diagonal}
 3(P_M'(t))^2-4P_M(t)P_M''(t)
 =48(M^2-1)(t^2-1)^2.
\end{equation}
If $M\geq1$, equations \eqref{eq:root} and \eqref{eq:diagonal} give concavity on the diagonal. If $0<M<1$, the second derivative is positive for every $t>0$ except $t=1$. This proves the first classification. It also excludes concavity on the full positive orthant when $0<M<1$. A concave function remains concave when restricted to an affine line.

To treat $M>1$, let
\[
 x_*=(2,2,1/2,1/2),\qquad v=(2,2,1,1),
 \qquad Q_M(s)=f_M(x_*+sv).
\]
This line lies in $\Gamma$ for $s>-1/2$. Direct expansion gives
\begin{align*}
 Q_M(s)={}&4s^4+(12M+12)s^3+(28M+26)s^2\\
          &+27(M+1)s+(41+40M)/4.
\end{align*}
Thus
\[
 Q_M(0)=\frac{41+40M}{4},\qquad
 Q_M'(0)=27(M+1),\qquad Q_M''(0)=52+56M.
\]
It follows that
\begin{equation}\label{eq:direction}
 4Q_M(0)Q_M''(0)-3(Q_M'(0))^2=(M-1)(53M+55)>0.
\end{equation}
By \eqref{eq:root}, the second derivative of $H_M$ along this line is positive at the interior point. Hence $H_M$ is not concave when $M>1$.

\section{Concavity when the parameter is one}
It remains to prove concavity when $M=1$. The subset expansion gives
\[
 f_1(x)=\prod_{i=1}^4(1+x_i).
\]
To prove concavity, use the following geometric-mean inequality.
For $u\in(0,\infty)^4$, put $G(u)=(u_1u_2u_3u_4)^{1/4}$. The arithmetic--geometric mean inequality implies, for positive $u,z$,
\begin{equation}\label{eq:amgm}
 G(u)\leq\frac{G(z)}4\sum_{i=1}^4\frac{u_i}{z_i}.
\end{equation}
For positive $u,w$, take $z=\theta u+(1-\theta)w$, where $0\leq\theta\leq1$. Apply \eqref{eq:amgm} to $u$ and $w$. Multiply each inequality by its corresponding weight and add. This gives
\[
 \theta G(u)+(1-\theta)G(w)
 \leq\frac{G(z)}4\sum_{i=1}^4
       \frac{\theta u_i+(1-\theta)w_i}{z_i}=G(z).
\]
Thus $G$ is concave. Composition with the affine map $x\mapsto(1+x_1,\ldots,1+x_4)$ proves concavity of $H_1$. Together with the preceding section, this proves Theorem~\ref{thm:classification}.

\section{A finite midpoint witness}
A separate rational certificate proves the failure at $M=2$ directly from the concavity inequality. Let
\[
 x_-=(3/2,3/2,1/4,1/4),\qquad
 x_+=(5/2,5/2,3/4,3/4).
\]
Their midpoint is $x_*$. Exact evaluation gives
\[
 f_2(x_-)=933/64,\quad f_2(x_*)=121/4,\quad f_2(x_+)=3597/64.
\]
The rational comparisons
\[
 (977/500)^4<933/64,\qquad (1369/500)^4<3597/64,
 \qquad121/4<(4691/2000)^4
\]
give
\[
 \frac{H_2(x_-)+H_2(x_+)}2
 >\frac{1173}{500}>\frac{4691}{2000}>H_2(x_*).
\]
This is the strict reverse of the required midpoint inequality.

\section{Formal verification and scope}
The Lean namespace \texttt{MathAndTheory.DiagonalConcavity} proves the polynomial identities, derivative formulas, and concavity statements. The declaration \path{complete_classification} proves Theorem~\ref{thm:classification}. The definition \texttt{elementary} uses Mathlib's \texttt{Multiset.esymm}. The theorem \path{family_eq_elementary} proves that this definition equals the coordinate formulas. The declaration \path{elementary_symmetric_counterexample} states the counterexample with these elementary symmetric functions. The declaration \path{midpoint_strict_violation} proves the midpoint certificate.

The result disproves Conjecture 1. It does not classify arbitrary coefficient vectors.

Lachaume's separate sufficient condition requires a real-rooted diagonal polynomial. The counterexample does not contradict that condition.

The Walsh-convolution note \cite{walsh} concerns Conjectures 4 and 5. It makes no claim for Conjectures 1 and 3. No claim of first priority is made.

\section*{Autonomous AI research}
This work is fully autonomous agentic research generated by AI systems. AI agents produced the mathematical arguments, formal proofs, computations, and manuscript. Dev D. Goyal is responsible for publication. This preprint has not been peer reviewed.

\bibliographystyle{plainnat}
\bibliography{references}
\end{document}
